% !TEX Root=Seminararbeit.tex

\section{Einleitung}
\onehalfspacing

Czesław Miłosz gehört zu den bedeutendsten polnischen Schriftstellern des 20. Jahrhunderts. Seine Bibliografie umfasst neben Lyrik auch eine Vielzahl von Essays, in denen er sich mit persönlichen, historischen und gesellschaftlichen Erfahrungen auseinandersetzt. Zu den persönlichen Erfahrungen, die in seinen Werken direkt oder auch indirekt Gegenstand literarischer Reflexion werden, gehören auch der Zweite Weltkrieg und die Okkupation Polens durch das Dritte Reich.
\\[6pt]
Diese beiden Ereignisse waren für die polnische Bevölkerung mit tiefgreifenden gesellschaftlichen und persönlichen Erfahrungen verbunden. Auch Czesław Miłosz erlebte diese Zeit in Polen und begann sich bereits während dieser Zeit in unterschiedlichen literarischen Formen mit seinen Erlebnissen und Erfahrungen auseinanderzusetzen.
\\[6pt]
Vor diesem Hintergrund untersucht diese Arbeit, wie Czesław Miłosz Erfahrungen von Krieg, Besatzung und moralischer Bedrohung literarisch verarbeitet. Im Zentrum steht dabei die Frage, mit welchen poetischen Mitteln Miłosz die deutsche Okkupation Polens reflektiert und wie sich darin individuelle Wahrnehmung und historische Erfahrung verbinden.
\\[6pt]
Als Untersuchungsgegenstand dient eine Auswahl lyrischer Texte von Czesław Miłosz. Im Zentrum stehen dabei Texte, die zwischen 1943 und 1945 entstanden sind. Ergänzend werden ein vor dem Zweiten Weltkrieg entstandener Text sowie zwei Texte aus der Nachkriegszeit untersucht. Durch diese zeitliche Erweiterung soll ein Vergleich ermöglicht werden, der Veränderungen in Miłoszs lyrischer Verarbeitung von Krieg und Okkupation sichtbar macht.